
\par 算法的时间和空间复杂度通常用大O记号表达。对于时间复杂度，计数的单位为算法执行过程中的基本运算（其工作量有常数上界）。在“算术模型”中，两个数的加、乘、除、比较等被视为单位运算；更精确的方式是“字节模型”，数值用二进制表示，按逐字节的运算计数。一般数值不会随算法运行而增大，使用算术模型即可；但考虑数值大小时，按字节计算是必要的（如素数判定问题）。

\begin{dfn} \textbf{大O记号}\\
设$f,g: \mathbb{N}\to \mathbb{R}^+$, 若存在$c>0$和正整数$N$，使得对任意$n\ge N$有$f(n)\le cg(n)$，则称$f(n)$是$O(g(n))$的。\index{大O记号}
\end{dfn}

\par 尽管文献中一般采用上述定义，事实上其中的条件等价于“存在$c>0$，使得对任意$n\ge 0$有$f(n)\le cg(n)$”。